\section{Grundlagen zeitkontinuierlicher Systeme}

Ein \textbf{zeitkontinuierliches System} beschreibt den Zusammenhang zwischen einer Eingangsgrösse $u(t)$ und einer Ausgangsgrösse $y(t)$
für eine kontinuierliche Zeitvariable $t \in \mathbb{R}$.
In der Leistungselektronik ist das typischerweise der Lastkreis (z.B.\ $R$-$L$, $R$-$C$, $R$-$L$-$C$), der durch \textbf{Differentialgleichungen}
modelliert wird.

\smallskip

\textbf{Zielbild (Outcome):}
\begin{itemize}
    \item Modell aufsetzen (DGL oder Zustandsraum)
    \item Zeitverlauf berechnen (Einschwingen, Ausschwingen, Umschalten)
    \item Kennwerte ableiten (Mittelwert, Effektivwert, Leistung, Ripple, Zeitkonstante)
    \item Stationären Betrieb bestimmen (periodisch, eingeschwungen)
\end{itemize}

\subsection{Systembegriffe und Klassifikation}

\subsubsection{Eingang, Ausgang, Zustand}
\begin{itemize}
    \item \textbf{Eingang} $u(t)$: Anregung von aussen (z.B.\ Quelle, PWM, Schalterstellung).
    \item \textbf{Ausgang} $y(t)$: beobachtete Grösse (z.B.\ Laststrom $i(t)$, Lastspannung $u_R(t)$).
    \item \textbf{Zustand} $x(t)$: minimaler Satz interner Variablen, der zusammen mit $u(t)$ die Zukunft vollständig bestimmt.
          In elektrischen Netzwerken sind Zustände typischerweise \textbf{Induktorströme} und \textbf{Kondensatorspannungen}.
\end{itemize}

\subsubsection{Linearität}
Ein System ist \textbf{linear}, wenn \textbf{Superposition} gilt:
\[
u_1(t)\rightarrow y_1(t),\; u_2(t)\rightarrow y_2(t)\quad \Rightarrow \quad
\alpha u_1(t)+\beta u_2(t)\rightarrow \alpha y_1(t)+\beta y_2(t).
\]
Lineare Modelle sind der Standard für $R$-$L$-$C$-Netze (bei idealen Bauteilen und ohne Sättigung).

\subsubsection{Zeitinvarianz}
Ein System ist \textbf{zeitinvariant}, wenn eine Zeitverschiebung am Eingang dieselbe Zeitverschiebung am Ausgang erzeugt:
\[
u(t-t_0)\rightarrow y(t-t_0).
\]
Ein \textbf{LTI-System} (linear time-invariant) ist der Kernfall für analytische Lösungen (Faltung, Laplace, Frequenzgang).

\subsubsection{Kausalität}
\textbf{Kausal} bedeutet: $y(t)$ hängt nur von $u(\tau)$ für $\tau \le t$ ab.
Physikalische elektrische Systeme sind kausal.

\subsubsection{Stabilität (praktisch relevant)}
\begin{itemize}
    \item \textbf{BIBO-stabil}: beschränkter Eingang $\Rightarrow$ beschränkter Ausgang.
    \item Für LTI-Systeme entspricht das (bei rationalen Übertragungsfunktionen) Polen mit negativer Realteil-Lage.
\end{itemize}

\subsection{Mathematische Modellierung}

\subsubsection{Differentialgleichungen als Systemmodell}
Viele zeitkontinuierliche Systeme werden durch eine lineare DGL beschrieben:
\[
a_n \frac{d^n y}{dt^n} + \dots + a_1 \frac{dy}{dt} + a_0 y
=
b_m \frac{d^m u}{dt^m} + \dots + b_1 \frac{du}{dt} + b_0 u.
\]
In der Leistungselektronik dominieren häufig Systeme \textbf{1.\ Ordnung} (z.B.\ $R$-$L$ oder $R$-$C$) und \textbf{2.\ Ordnung} (z.B.\ $L$-$C$ mit Dämpfung).

\subsubsection{Zustandsraumdarstellung}
Alternative Standardform (skalierbar, schaltertauglich):
\[
\dot{x}(t)=A x(t)+B u(t), \qquad y(t)=C x(t)+D u(t).
\]
\textbf{Use-Case:} Bei Schaltvorgängen ändert sich typischerweise $A,B,C,D$ stückweise (Topologie 1, Topologie 2, \dots).
Das ist der saubere Rahmen für \textbf{stückweise lineare} Systeme.

\subsection{Systeme 1.\ Ordnung (Business Case: $R$-$L$ und $R$-$C$)}

\subsubsection{$R$-$L$-System}
Grundgleichung (typisch in Übungen):
\[
u(t)=R\,i(t)+L\,\frac{di(t)}{dt}.
\]
Zeitkonstante:
\[
\cbl{\tau=\frac{L}{R}}
\]
Interpretation: Nach ca.\ $t=5\tau$ ist ein Einschwingvorgang praktisch abgeschlossen (ca.\ 99\%).

\subsubsection{$R$-$C$-System}
Grundgleichung:
\[
i(t)=C\frac{du_C(t)}{dt},\qquad u(t)=u_R(t)+u_C(t)=R i(t)+u_C(t).
\]
Zeitkonstante:
\[
\cbl{\tau=R C}
\]

\subsection{Lösungskonzept für lineare DGL 1.\ Ordnung}

Standardform:
\[
\frac{dy}{dt}+a\,y=b\,u(t).
\]

\subsubsection{Homogene Lösung (natürliche Antwort)}
Setze $u(t)=0$:
\[
\frac{dy_h}{dt}+a\,y_h=0 \quad \Rightarrow \quad y_h(t)=K\,e^{-a t}.
\]
\textbf{Interpretation:} Das System trägt Energie ab. Exponentielles Abklingen.

\subsubsection{Partikuläre Lösung (erzwungene Antwort)}
Für einen konstanten Eingang $u(t)=U_0$:
\[
\frac{dy_p}{dt}+a\,y_p=b\,U_0.
\]
Typischer Ansatz: $y_p=\text{konstant}$.
Dann ergibt sich der Endwert:
\[
y_\infty=\frac{b}{a}U_0.
\]

\subsubsection{Gesamtlösung}
\[
y(t)=y_p(t)+y_h(t).
\]
Für den konstanten Eingang ist das die klassische Einschwingform:
\[
\cbl{y(t)=y_\infty + \big(y(0)-y_\infty\big)\,e^{-t/\tau}}
\]
Hier gilt $\tau=1/a$.

\subsubsection{Initialbedingungen und Kontinuität}
\begin{itemize}
    \item \textbf{Induktorstrom ist stetig:} $\cbl{i_L(t^-)=i_L(t^+)}$ (Strom kann nicht sprunghaft ändern, sonst bräuchte man unendliche Spannung).
    \item \textbf{Kondensatorspannung ist stetig:} $\cbl{u_C(t^-)=u_C(t^+)}$ (Spannung kann nicht sprunghaft ändern, sonst bräuchte man unendlichen Strom).
\end{itemize}
Diese Kontinuitätsregeln sind bei \textbf{Schaltvorgängen} die zentrale Randbedingung zur Bestimmung der Integrationskonstanten.

\subsection{Stückweise Systeme bei Schaltbetrieb (Topologie-Wechsel)}

In Schaltungen mit Schaltern gilt oft:
\[
\text{Topologie 1: } \dot{x}=A_1 x + B_1 u,\quad t\in[t_0,t_1)
\]
\[
\text{Topologie 2: } \dot{x}=A_2 x + B_2 u,\quad t\in[t_1,t_2)
\]
\textbf{Vorgehen:}
\begin{enumerate}
    \item Pro Zeitintervall DGL lösen (typisch exponentiell).
    \item Übergangsbedingungen am Schaltzeitpunkt mit Kontinuität ansetzen.
    \item Bei periodischem Betrieb zusätzlich Periodizität verlangen:
    \[
    x(t_0)=x(t_0+T_s)
    \]
    \item Daraus Integrationskonstanten bestimmen.
\end{enumerate}

\subsection{Stationärer Betrieb vs.\ Transient}

\subsubsection{Transient (Einschwingen)}
Startbedingungen sind relevant. Lösung enthält $e^{-t/\tau}$-Anteile, die mit der Zeit verschwinden.

\subsubsection{Stationär (eingeschwungen)}

\crd{\textrightarrow\ Stationär bedeutet periodisch wiederholbar, nicht konstant.}

Die natürlichen Anteile sind ``abgebaut''. Im periodischen Schaltbetrieb bedeutet stationär:
\[
\cbl{x(t)=x(t+T_s)}
\]
Wichtig: Stationär heisst nicht konstant, sondern \textbf{periodisch wiederholbar}.

\subsection{Übertragungsfunktion und Laplace }

\subsubsection{Laplace-Transformation}
Für LTI-Systeme ist Laplace ein Standard-Werkzeug zur Formalisierung von Dynamik:
\[
\cbl{\mathcal{L}\{f(t)\}=F(s)=\int_0^\infty f(t)\,e^{-st}\,dt}
\]

\subsubsection{Übertragungsfunktion}
Bei Nullanfangsbedingungen:
\[
\cbl{G(s)=\frac{Y(s)}{U(s)}}
\]
Für $R$-$L$ mit Ausgang $i(t)$ und Eingang $u(t)$:
\[
\cbl{G(s)=\frac{I(s)}{U(s)}=\frac{1}{R+Ls}}
\]
Pol:
\[
\cbl{s=-\frac{R}{L}}
\]
\textbf{Interpretation:} Negative Pol-Lage $\Rightarrow$ stabil, exponentielles Abklingen.

\subsubsection{Impulsantwort und Faltung}
Für LTI gilt:
\[
y(t)=(g*u)(t)=\int_{0}^{t} g(\tau)\,u(t-\tau)\,d\tau,
\]
wobei $g(t)$ die Impulsantwort ist.
In der Leistungselektronik wird das oft indirekt genutzt (Filterwirkung, Glättung, Spektralargumentation).

\subsection{Kennwerte im Zeitbereich (Operational KPIs)}

\subsubsection{Mittelwert (DC-Anteil)}
Für periodische Signale:
\[
\cbl{\overline{x}=\frac{1}{T}\int_{0}^{T} x(t)\,dt}
\]

\subsubsection{Effektivwert (RMS)}
\[
\cbl{X_{\mathrm{RMS}}=\sqrt{\frac{1}{T}\int_{0}^{T} x^2(t)\,dt}}
\]
\textbf{Use-Case:} Thermische Bewertung und Leistungsberechnung.

\subsubsection{Wirkleistung}
\[
\cbl{p(t)=u(t)\,i(t),\qquad P=\overline{p(t)}=\frac{1}{T}\int_0^T u(t)i(t)\,dt}
\]
Für ohmsche Last:
\[
P=R\,I_{\mathrm{RMS}}^2=\frac{U_{\mathrm{RMS}}^2}{R}.
\]

\subsection{Typische Fehlerbilder (Quality Gate)}

\crd{Diese Punkte sind typische Prüfungsfallen.}

\begin{itemize}
    \item Kontinuität falsch angesetzt (Induktorstrom oder Kondensatorspannung springt).
    \item Stationär mit konstant verwechselt (periodisch vs.\ DC).
    \item Zeitvariable im Exponenten falsch skaliert (fehlende $\tau$, falsches Vorzeichen).
    \item Periodenbedingungen nicht geschlossen (Integrationskonstanten bleiben frei).
    \item RMS und Mittelwert vermischt (insbesondere bei nichtsinusförmigen Signalen).
\end{itemize}

\subsection{Minimal-Checkliste für Aufgaben (Execution Plan)}
\begin{enumerate}
    \item Topologie pro Zeitintervall identifizieren.
    \item DGL aufstellen (KVL/KCL).
    \item Standardform erkennen ($1.\ Ordnung$: exponentiell).
    \item Anfangs- und Übergangsbedingungen definieren (Kontinuität).
    \item Bei stationär-periodisch: $x(t)=x(t+T_s)$ einsetzen.
    \item Gesuchte Grössen ableiten: $u_R=Ri$, $u_L=L\,di/dt$, Leistung, RMS, Mittelwert.
\end{enumerate}
